% Template:     Tesis LaTeX
% Documento:    Archivo de ejemplo
% Versión:      3.4.2 (07/02/2025)
% Codificación: UTF-8
%
% Autor: Pablo Pizarro R.
%        pablo@ppizarror.com
%
% Manual template: [https://latex.ppizarror.com/tesis]
% Licencia MIT:    [https://opensource.org/licenses/MIT]

% ------------------------------------------------------------------------------
% NUEVO CAPÍTULO
% ------------------------------------------------------------------------------
% A diferencia de Template-Informe, Template-Tesis requiere el uso de capítulos; las secciones, subsecciones, etc son parte de un capítulo. Se recomienda el uso de un capítulo en un archivo distinto
\chapter{Introduction}

\section{Clinical and technical motivation}
\section{Problem definition: brain tumor classification in MRI}
\section{Limitations of conventional discriminative models/ explaining discriminative and discriminative models}
\section{Proposed solution: generative inference + explainability}
\section{General and specific objectives}
\section{Hypothesis}


% ------------------------------------------------------------------------------
% NUEVO CAPÍTULO
% ------------------------------------------------------------------------------
\chapter{Background}

\section{Basic neuroanatomy and brain tumor types}
\section{Magnetic Resonance Imaging (MRI): principles and challenges}
\section{Clinical diagnostic processes}
\section{Role of artificial intelligence in medical imaging}

	



% ------------------------------------------------------------------------------
% NUEVO CAPÍTULO
% ------------------------------------------------------------------------------
\chapter{Literature Review}

\section{Discriminative models (CNN, Transfer Learning)}

\subsection{Comparative table: discriminative classifiers in MRI}

\section{Generative models (GAN, CGAN, VAE, Diffusion Models, Normalizing Flows)}

\subsection{Comparative table: generative classifiers in MRI
}
\section{Classification by generation (Inference by Generation)}

\section{Explainable AI (XAI) in medical imaging (Grad-CAM, LIME, SHAP,llm)}

\section{SR methods}
\section{human in the loop}
\section{Existing comparative studies and research gaps}

\section{Justification for the proposed approach}
\
% ------------------------------------------------------------------------------
% REFERENCIAS, revisar configuración \stylecitereferences
% ------------------------------------------------------------------------------
\bibliography{library}


% ------------------------------------------------------------------------------
% ANEXO
% Existe adicionalmente el entorno \begin{appendixd} que permite insertar
% \chapter y el entorno \begin{appendixdtitle}[style1] (4 estilos diferentes),
% el cual acepta \chapter y escribe el título de anexos encima
% ------------------------------------------------------------------------------
\begin{appendixs}
	
	\section{Cálculos realizados}

	\subsection{Metodología}
	\lipsum[1-2]

	% Imagen, se numerará automáticamente con la letra del anexo según
	% la configuración \appendixindepobjnum
	\insertimage[\label{img:anexo-2}]{ejemplos/test-image.png}{scale=0.25}{Imagen en anexo.}

	\subsection{Resultados}
	\lipsum[10]

	% Tablas
	\enabletablerowcolor[2] % Activa el color de celda
	\begin{table}[H]
		\begin{threeparttable}
		\centering
		\caption{Tabla de cálculo.}
		\begin{tabular}{cccC{4cm}}
			\hline
			\textbf{Elemento} & $\epsilon_i$ & \textbf{Valor} & \textbf{Descripción} \bigstrut \\
			\hline
			A     & 10    & 3,14$\pi$ & Valor muy interesante\tnote{a} \\
			B     & 20    & 6 & Segundo elemento \\
			C     & 30    & 7 & Tercer elemento\tnote{1} \\
			D     & 150    & 10 & Sin descripción \\
			E     & 0    & 0 & Cero \\
			\hline
			\end{tabular}
		\begin{tablenotes}
			\item[a] Este elemento tiene una descripción debajo de la tabla
			\item[1] Más comentarios
		\end{tablenotes}
		\end{threeparttable}
		\label{tab:anexo-1}
	\end{table}
	\disabletablerowcolor % Desactiva el color de celda

\end{appendixs}